\scanpage{401}\mspage{315}

Les compatibilités $(****)$ impliquent qu'on a des hom. \deleted{\unsure{de}}

\diagram{\begin{tikzcd}[ampersand replacement=\&,column sep=large]
\unsure{\varphi|\phi}(i)\arrow[r,shift left=0.65ex,"u"] \& j\arrow[l,shift left=0.65ex,"v"] \& {,\ \text{où }\ j\overset{\text{\unsure{def|déf}}}{=}\underset{\alpha\in I_{\alpha_\pi}}{\underset{\longrightarrow}{\lim}}\,j_\alpha}
\end{tikzcd}}

\inserted{(\unsure{existe|existant} \deleted{\unsure{la}}\unsure{déjà|la|là} \uline{dans $J'$} donc dans $J$\unsure{)|)'}},

et les compatibilités $(**)$ impliquent que

\[
\deleted{\unsure{u}\ \text{et}\ \unsure{v}\ \text{\unsure{son}}}\qquad
vu=\mathrm{id}_{\varphi(i)}\deleted{\unsure{\varphi(i)}},\quad vu=\mathrm{id}_j,
\]

donc que $u,v$ sont des iso. inverses l'un de l'autre. \deleted{Or la propriété 4) pour $J$ implique que $j\in J'$,} Comme $J'$ est une sous-catégorie \uline{str.} pleine de $J$, il s'ensuit bien que $\varphi(i)\in J'$, i.e. $i\in I'$, qed.

\deleted{\unsure{Donc}} \unsure{Reste|Rest} \unsure{maint$^{\mathrm{t}}$|maintenant} à construire les

\[
\underset{\mathrm{Ob}\,I}{\underset{\unsure{\mathord{\tikz[baseline=(c.base)]{\node[rotate=-90,inner sep=0pt] (c) {$\scriptstyle\in$};}}}}{i_\alpha}},\tau_{\beta\alpha},\quad
\underset{\mathrm{Ob}\,J'}{\underset{\unsure{\mathord{\tikz[baseline=(c.base)]{\node[rotate=-90,inner sep=0pt] (c) {$\scriptstyle\in$};}}}}{j_\alpha}},\tau'_{\beta\alpha},\quad u_\alpha,v_\alpha
\]

satisfaisant aux contraintes \deleted{$\unsure{(*)}$} $(**)\,(***)\,(****)$. On va faire la construction en trois étapes.

\textcircled{1} Départ :

\[
i_0=\unsure{0|i},
\]

Comme $J'$ est cofinal dans $J$, existe $j_0\in\mathrm{Ob}\,J'$ et $\varphi(i_0)\xrightarrow{u_0}j_0$.

% archive sheet number: 400
